%%═══════════════════════════════════════════════════════════════════
%% Lecture 04 — Deuteron Ground State & Nuclear Shell Model
%% 77603 — Nuclear Physics
%% Date: 5 May 2026
%% Lecturer: Dr. Moshe Friedman
%%═══════════════════════════════════════════════════════════════════

\section{Lecture 04 — Deuteron Ground State and the Nuclear Shell Model}
\label{sec:lec04}
\date{5 May 2026}

%%───────────────────────────────────────────────────────────────────
\subsection*{Overview}
%%─────────────────────
This lecture has two parts.  The first solves the quantum mechanics of the
deuteron: we determine its ground-state angular momentum by combining parity
and magnetic-moment arguments, discover that the ground state is a
superposition of $l=0$ and $l=2$ components (proving the nuclear force has a
tensor part), and solve the Schrödinger equation for a square-well potential to
extract the potential depth $V_0\approx35\,\mathrm{MeV}$.  The second part
introduces the nuclear shell model: we motivate the existence of nuclear magic
numbers, justify why independent-particle motion is not unreasonable, try
several model potentials and fail, and finally show that a spin-orbit coupling
term $V_\mathrm{so}\propto\vec{l}\cdot\vec{s}$ reorganizes the levels so as
to reproduce all observed magic numbers.


%%───────────────────────────────────────────────────────────────────
\section{Part I — Solving the Deuteron}
%%───────────────────────────────────────────────────────────────────

\subsection{Determining the Orbital Angular Momentum}

\subsubsection*{Physical Motivation}
We know from Lecture~03 that the deuteron binding energy is
$B_d\approx2.22\,\mathrm{MeV}$ and that its spin-parity is $I^\pi=1^+$.  We
now ask: what orbital angular momentum $l$ does the proton–neutron pair carry?
The answer is not trivial, and it turns out to be more interesting than a simple
integer.

The measured spin-parity $I^\pi=1^+$ constrains the possible $l$ values.
Parity of the spatial part of the wavefunction is $(-1)^l$.  For the nuclear
state to have $\pi=+1$ we need $(-1)^l=+1$, so $l$ must be \emph{even}: $l=0,
2, 4, \ldots$.  Including the spin of the nucleon pair $S=0$ (singlet) or $S=1$
(triplet) and requiring the total $J=1$ reduces the candidates further:
\begin{itemize}
  \item $l=0,\;S=1$: $J=1$, $\pi=+1$\quad$\checkmark$
  \item $l=2,\;S=1$: $J=1,2,3$ (all have $\pi=+1$), $J=1$ is possible\quad$\checkmark$
  \item $l=1,\;S=0$: $J=1$, $\pi=-1$\quad$\times$ (wrong parity)
  \item $l=1,\;S=1$: $J=0,1,2$, $\pi=-1$\quad$\times$ (wrong parity)
\end{itemize}
So \textbf{two choices survive: $l=0$ or $l=2$ (or a mixture)}.

\subsubsection*{Using the Magnetic Moment to Distinguish}

\paragraph{Prediction for $l=0$ (pure S-wave).}
If $l=0$ there is no orbital contribution and the magnetic moment comes entirely
from the spins.  With both nucleons in the triplet $S=1$ state:
\begin{equation}
  \label{eq:lec04:mu_Swave}
  \mu_{l=0} = \frac{1}{2}\bigl[g_s(p)+g_s(n)\bigr]\,\mu_N
  = \frac{1}{2}(5.586-3.826)\,\mu_N
  \approx 0.879\,\mu_N.
\end{equation}
The measured value is
\begin{equation}
  \label{eq:lec04:mu_measured}
  \mu_d = 0.8574\,\mu_N.
\end{equation}
These are close but do \emph{not} agree — the discrepancy is well outside
experimental uncertainty.

\paragraph{Prediction for $l=2$ (pure D-wave).}
A more involved calculation gives $\mu_{l=2}\approx0.31\,\mu_N$, which is even
further from the measured value.

\paragraph{Conclusion: a mixed S–D ground state.}
Neither pure $l=0$ nor pure $l=2$ fits.  The resolution is that the deuteron
ground state is a \emph{superposition}:
\begin{equation}
  \label{eq:lec04:psi_d}
  \psi_d = A\,\psi_{S\text{-wave}}\;+\;B\,\psi_{D\text{-wave}},
\end{equation}
with normalization $|A|^2+|B|^2=1$.  Matching the measured $\mu_d$ fixes the
D-wave probability:
\begin{equation}
  \label{eq:lec04:d_wave}
  \boxed{|B|^2 \equiv P_D \approx 0.04.}
\end{equation}
\textit{Physical significance:} About $4\%$ of the time the proton and neutron
are in an $l=2$ state.  Because $l$ is not fixed, $l$ is \emph{not a good
quantum number} for the deuteron.  A central potential cannot mix $l=0$ and
$l=2$; therefore the nuclear force contains a \textbf{non-central (tensor)
component}.  The tensor force is the dominant deviation from a simple central
interaction.

\subsection{Square-Well Model for the Deuteron}

\subsubsection*{Physical Motivation}
With $P_D\approx4\%$ being small, we can work to good approximation in the
$l=0$ sector alone.  The goal is to relate the measured binding energy
$B_d\approx2.22\,\mathrm{MeV}$ and the known nuclear radius $R\approx2\,\mathrm{fm}$
to the depth $V_0$ of the nuclear potential.  The simplest possible potential that
is short-ranged and attractive is the square well.

\paragraph{Reduced-mass setup.}
Working in the centre-of-mass frame, the two-body problem reduces to a
one-body problem with the reduced mass
\begin{equation}
  \label{eq:lec04:reduced_mass}
  \mu = \frac{m_p m_n}{m_p+m_n}\approx\frac{m_N}{2},
\end{equation}
and the relative coordinate $\vec{r}$.  The wavefunction factorises as
$\psi(\vec{r})=R_{nl}(r)Y_l^m(\theta,\phi)$.  Setting $u(r)=r R(r)$ and
$l=0$, the radial Schrödinger equation becomes
\begin{equation}
  \label{eq:lec04:radial_eq}
  -\frac{\hbar^2}{2\mu}\dv[2]{u}{r}+V(r)u(r)=Eu(r).
\end{equation}
For a bound state $E=-B<0$.

\paragraph{Square-well potential.}
\begin{equation}
  \label{eq:lec04:square_well}
  V(r)=\begin{cases}-V_0&r<R,\\0&r\geq R.\end{cases}
\end{equation}

\paragraph{Solutions in each region.}
Define
\begin{equation}
  k_1=\sqrt{\frac{2\mu(V_0-B)}{\hbar^2}},\qquad
  k_2=\sqrt{\frac{2\mu B}{\hbar^2}}.
\end{equation}
\begin{itemize}
  \item \textbf{Inside} ($r<R$): equation $u''+k_1^2 u=0$; general solution
    $A\sin k_1 r+C\cos k_1 r$.  The boundary condition $u(0)=0$ (so that
    $\psi(0)=u(0)/r$ stays finite) forces $C=0$:
    \begin{equation}
      u_\mathrm{in}(r)=A\sin k_1 r.
    \end{equation}
  \item \textbf{Outside} ($r>R$): equation $u''-k_2^2 u=0$; general solution
    $De^{-k_2 r}+Fe^{+k_2 r}$.  Normalizability forces $F=0$:
    \begin{equation}
      u_\mathrm{out}(r)=De^{-k_2 r}.
    \end{equation}
\end{itemize}

\paragraph{Matching boundary conditions at $r=R$.}
Continuity of $u$ and $u'$:
\begin{align}
  A\sin k_1 R &= De^{-k_2 R},\\
  Ak_1\cos k_1 R &= -k_2 De^{-k_2 R}.
\end{align}
Dividing eliminates $A$ and $D$:
\begin{equation}
  \label{eq:lec04:bc}
  \boxed{k_1\cot(k_1 R)=-k_2.}
\end{equation}
\textit{Physical significance:} This transcendental equation, together with the
definitions of $k_1$ and $k_2$ (which depend on $V_0$, $B$, $\mu$), uniquely
determines $V_0$ once $B_d$ and $R$ are known from experiment.

\paragraph{Numerical result.}
Inserting $B_d=2.22\,\mathrm{MeV}$, $R\approx2\,\mathrm{fm}$, $\mu\approx
m_N/2$:
\begin{equation}
  \label{eq:lec04:V0}
  \boxed{V_0\approx35\,\mathrm{MeV}.}
\end{equation}
The deuteron sits near the very edge of the well: $k_1 R\approx\pi/2$,
meaning the wavefunction inside is just reaching its first quarter-oscillation
when it meets the wall.

\subsubsection*{Minimum Potential Depth for a Bound State}
A bound state requires $k_1 R\geq\pi/2$, i.e.\
\begin{equation}
  \label{eq:lec04:Vmin}
  V_{0,\min}=\frac{\pi^2\hbar^2}{8\mu R^2}\approx25\,\mathrm{MeV}.
\end{equation}
\textit{Physical significance:} The actual depth $V_0\approx35\,\mathrm{MeV}$ is only
$40\%$ above the minimum needed.  The deuteron is barely bound.  Had the nuclear
force been slightly weaker, no bound state of the two nucleons would exist, and
nuclear chemistry as we know it would be impossible.


%%───────────────────────────────────────────────────────────────────
\section{Part II — The Nuclear Shell Model}
%%───────────────────────────────────────────────────────────────────

\subsection{Magic Numbers and Empirical Evidence}

\subsubsection*{Physical Motivation}
Atoms are stable when electron shells are complete (noble gases).  By analogy,
nuclei might be especially stable when their proton or neutron numbers fill
particular ``shells.''  The numbers at which this happens are called
\textbf{magic numbers}:
\begin{equation}
  \label{eq:lec04:magic}
  \boxed{2,\quad 8,\quad 20,\quad 28,\quad 50,\quad 82,\quad 126.}
\end{equation}
Several independent experimental observations converge on the same numbers:
\begin{enumerate}
  \item \textbf{Number of stable isotopes/isotones} — magic-$N$ or magic-$Z$
    nuclides have unusually many stable species.  For example, $N=28$ has
    5--6 stable isotones; neighbouring $N$ values have only 1--2.
  \item \textbf{Alpha-decay energies} — the kinetic energy released in $\alpha$
    decay of nuclei with $N=126$ (a filled neutron shell) drops sharply as
    $N$ crosses 126, then rises again.
  \item \textbf{Neutron-capture cross sections} — nuclei at magic $N$ have
    anomalously \emph{small} cross sections: the full shell resists adding one
    more neutron.
  \item \textbf{Nuclear radii} — the empirical $R=r_0 A^{1/3}$ trend shows
    kinks at magic numbers.
\end{enumerate}
The pattern unmistakably echoes atomic shell closure.

\subsection{Why Independent-Particle Motion is Plausible}

\subsubsection*{Physical Motivation}
At first sight, treating nucleons as free particles in a mean potential seems
absurd: the nucleus is a dense system where every nucleon interacts strongly
with its neighbours.  Yet the Pauli exclusion principle saves the picture.

In the ground state, all orbitals up to the Fermi energy are occupied.  For two
nucleons to scatter, one must transfer enough energy to the other to kick it
into an unoccupied orbital above the Fermi level.  But \emph{all nearby states
are already full}.  The energy required is $\sim E_F\approx33\,\mathrm{MeV}$,
far larger than the thermal energy scale.  As a result, nucleons effectively
\emph{pass through each other without scattering}, and each nucleon moves in
a slowly varying mean-field potential $V(r)$ created by all the others.  This
is the \textbf{independent-particle approximation}.

\subsection{Mean-Field Potentials and their Failures}

\paragraph{Infinite square well.}
Solving $-\frac{\hbar^2}{2m}\nabla^2\psi=E\psi$ inside a sphere of radius $R$
gives energy levels labelled by $(n_r,l)$.  Filling them with $2(2l+1)$
nucleons each (factor 2 for spin) produces cumulative occupations that hit
shell gaps at
\[
  2,\;8,\;20,\;40,\;58,\;92,\;\ldots
\]
The first three (2, 8, 20) agree with experiment; beyond 20 the model fails.

\paragraph{Harmonic oscillator.}
Energy levels $E_N=\hbar\omega(N+3/2)$ with degeneracy $(N+1)(N+2)/2$ per
oscillator quantum $N=2n_r+l$.  Shell closures occur at
\[
  2,\;8,\;20,\;40,\;70,\;\ldots
\]
Again, only the first three magic numbers are reproduced.

\paragraph{Woods–Saxon potential.}
A more realistic rounded well,
\begin{equation}
  \label{eq:lec04:woods_saxon}
  \boxed{V(r)=\frac{-V_0}{1+\exp\!\left(\frac{r-R}{a}\right)},}
\end{equation}
where $R\approx r_0 A^{1/3}$ is the nuclear radius and $a\approx0.5\,\mathrm{fm}$
is the skin thickness, gives a softer level sequence.  It agrees better than
the square well with observed radii and separation energies, but still fails to
reproduce magic numbers beyond 20.

\textit{Physical significance:} The parameter $a$ in the Woods–Saxon potential
is the same skin thickness that appears in the charge-density profiles measured
by electron scattering in Lecture~02, confirming self-consistency.

\subsection{Spin-Orbit Coupling Solves the Problem}

\subsubsection*{Physical Motivation}
In atomic physics, the electron's spin couples to the magnetic field generated
by its orbital motion, splitting each $l$-level into $j=l+\frac{1}{2}$ and
$j=l-\frac{1}{2}$ sub-levels.  Maria Goeppert Mayer and J.\ Hans D.\ Jensen
showed (1949, Nobel Prize 1963) that adding a spin-orbit term to the nuclear
mean field,
\begin{equation}
  \label{eq:lec04:Vso}
  V_\mathrm{so}(r) = -V_\mathrm{ls}\,f(r)\,\vec{l}\cdot\vec{s},
\end{equation}
with $V_\mathrm{ls}>0$ and $f(r)\propto-dV_\mathrm{WS}/dr$, generates the
correct magic numbers.  The sign is important: in \emph{nuclei} the $j=l+\frac{1}{2}$
state is pushed \emph{down} (more bound), opposite to the atomic case.

\paragraph{Evaluating $\vec{l}\cdot\vec{s}$.}
Using $\vec{j}=\vec{l}+\vec{s}$ and squaring:
\begin{equation}
  \label{eq:lec04:LS_identity}
  \vec{l}\cdot\vec{s}=\frac{1}{2}\bigl(\vec{j}^2-\vec{l}^2-\vec{s}^2\bigr).
\end{equation}
In a state $|j,l,s\rangle$ with $s=\frac{1}{2}$:
\begin{equation}
  \label{eq:lec04:LS_expectation}
  \boxed{
  \langle\vec{l}\cdot\vec{s}\rangle =
  \begin{cases}
    +\dfrac{\hbar^2\,l}{2} & j=l+\tfrac{1}{2},\\[6pt]
    -\dfrac{\hbar^2\,(l+1)}{2} & j=l-\tfrac{1}{2}.
  \end{cases}
  }
\end{equation}
The energy splitting between the two sub-levels is
\begin{equation}
  \label{eq:lec04:LS_splitting}
  \Delta E = V_\mathrm{ls}\,f_0\,\frac{\hbar^2}{2}\,(2l+1),
\end{equation}
which \emph{grows with $l$}.  For large $l$ (e.g.\ $l=4$, the $g$-shell), the
splitting is large enough to push a level from one major oscillator shell into
the one below it, rearranging the cumulative level occupations and producing the
correct gaps at 28, 50, 82, and 126.

\textit{Physical significance:} Before Mayer and Jensen, the failure at 28, 50,
82, 126 was an embarrassment.  Adding a single physically motivated term — the
spin-orbit interaction — solved the entire problem at once.

\subsection{The Extreme Independent-Particle Model (EIPM)}

\subsubsection*{Physical Motivation}
Once the shell structure is established, predicting the ground-state spin and
parity of a nucleus becomes a counting exercise.  For simplicity we use the
\textbf{Extreme Independent-Particle Model}: even-$A$ nuclei have $I=0$ (all
pairs cancel), and for odd-$A$ nuclei \emph{only the last unpaired nucleon
determines $I^\pi$}.

\paragraph{Rule.}
Fill the proton and neutron shells separately up to the $A$-th nucleon.  The
last unpaired nucleon occupies a state with orbital quantum number $l$ and
total angular momentum $j$.  Then:
\begin{equation}
  \label{eq:lec04:EIPM}
  \boxed{I^\pi = j^\pi,\qquad \pi=(-1)^l.}
\end{equation}

\paragraph{Example: $\nucleus{O}{17}$ ($Z=8$, $N=9$).}
The eight protons fill $1s_{1/2}$, $1p_{3/2}$, $1p_{1/2}$ completely; since
$Z=8$ is magic, protons contribute $I=0$.  The nine neutrons fill the same
sequence plus one extra neutron that must go into the $1d_{5/2}$ level:
\begin{center}
\begin{tabular}{lcc}
\toprule
Level & Capacity & Cumulative $N$ \\
\midrule
$1s_{1/2}$ & 2 & 2 \\
$1p_{3/2}$ & 4 & 6 \\
$1p_{1/2}$ & 2 & 8 \\
$1d_{5/2}$ & 6 & 14 \\ \bottomrule
\end{tabular}
\end{center}
The 9th neutron sits in $1d_{5/2}$, so $j=\frac{5}{2}$, $l=2$, $\pi=+1$:
\begin{equation}
  I^\pi(\nucleus{O}{17}) = \tfrac{5}{2}^+.\quad \checkmark
\end{equation}
(The experimental value is indeed $\frac{5}{2}^+$.)

\paragraph{Example: $\nucleus{O}{15}$ ($Z=8$, $N=7$).}
Seven neutrons fill $1s_{1/2}$ (2), $1p_{3/2}$ (4), and one of the $1p_{1/2}$
states.  The last neutron is in $1p_{1/2}$: $j=\frac{1}{2}$, $l=1$,
$\pi=-1$:
\begin{equation}
  I^\pi(\nucleus{O}{15}) = \tfrac{1}{2}^-.\quad \checkmark
\end{equation}

\subsubsection*{Successes and Limitations}
The EIPM reproduces the \emph{spin-parity} of most odd-$A$ ground states correctly.
Magnetic moments are only reproduced qualitatively (the predicted values lie
between the Schmidt lines; real nuclei cluster near the lines but are not on
them).  Electric quadrupole moments are sometimes off by large factors, signalling
that collective deformation beyond the single-particle picture is present.


%%───────────────────────────────────────────────────────────────────
\subsection*{Summary of Key Results}
%%─────────────────────

\begin{enumerate}
  \item Deuteron ground state: $\psi_d=A\psi_S+B\psi_D$,
    $P_D=|B|^2\approx4\%$.  The $l=0$ and $l=2$ admixture proves the nuclear
    force has a tensor component.
  \item Square-well model: boundary condition $k_1\cot(k_1 R)=-k_2$ gives
    $V_0\approx35\,\mathrm{MeV}$; minimum for a bound state is
    $V_{0,\min}\approx25\,\mathrm{MeV}$.
  \item Magic numbers 2, 8, 20 follow from both the infinite-square-well and
    the harmonic-oscillator potential.  Reproducing 28, 50, 82, 126 requires
    spin-orbit coupling $V_\mathrm{so}\propto-\vec{l}\cdot\vec{s}$.
  \item EIPM prediction: for odd-$A$, $I^\pi=j^\pi$ of the last unpaired
    nucleon, with $\pi=(-1)^l$.
\end{enumerate}
